\documentclass[twocolumn]{article}
\usepackage[utf8]{inputenc}
\usepackage{abstract}
\usepackage{graphicx}
\usepackage{amsmath, amssymb}
\usepackage{hyperref}
\usepackage{caption}
\usepackage{subcaption}
\usepackage{geometry}
\usepackage{booktabs}
\usepackage{float}
\usepackage[backend=biber, style=apa]{biblatex}
\DeclareLanguageMapping{english}{english-apa}
\addbibresource{~/Documents/Zotero/references.bib}
\geometry{margin=1in}

\title{Using Logistic Regression and XGBoost for Bot Detection}
\author{
 Elias Schlie \\
 u783766 \\
 Tilburg University \\
  \texttt{e.p.schlie@tilburguniversity.edu}
}
\date{}

\begin{document}


\twocolumn[
  \maketitle

  \begin{abstract}
 The activity of bots on social media platforms is increasingly becoming a problem and platform providers are seeking ways to prevent bots from infiltrating their systems.
The following project evaluates two simple machine learning algorithms on their ability to discern between bots and users based on their usage data. 
The algorithms applied are logistic regression, which was able to reach an accuracy of $97\%$, and XGBoost, \cite{chenXGBoostScalableTree2016} which reached $98\%$ accuracy.
  \end{abstract}
  \vspace{3em}
]

\section{Introduction}

With the rise of large language models, bots now have the ability to immitate human behavior better than ever before. This makes descerning them from actual users increasingly hard, and more sophisticated methods are needed to detect them.

% Describe models

\section{Data and Preprocessing}
% Describe the dataset, relevant features, and how you prepared/split the data.

\subsection{Dataset Description}

The 
"bots\_vs\_users.csv" dataset contains information on $5874$ social media accounts, exactly half of which are bots. Every account has $59$ features, that include $14$ numerical ones like their \texttt{average\_views} and $45$ categorical ones like the city they live in. $41$ of the categorical features are binary like the the variable \texttt{has\_profile\_picture} for example.
Most of the features include a large number of missing or unknown values. 4483 accounts don't contain any numerical information, while for the individual categorical features, between 5006 and 24 of the accounts cary the label "unknown".
Only the target varialbe "target", which indicates if the account is a bot or not has no missing features.

\subsection{Data split}

Before any preprocessing was applied, the data was split into a training and a test set using a stratified 80/20 split. This was done prior to preprocessing due to the reason that some of the preprocessing steps learn from the data and thus would lead to data leakage if applied to the entire dataset at once. Model hyperparameters were validated based on a k-fold cross validation on the training set. To ensure reproducibility, this and every following step that involves randomness was seeded with the seed $42$.

\subsection{Preprocessing and Feature Engineering}

All the preprocessing and feature engineering steps were implemented in a single \texttt{DataCleaner} class. The class's \texttt{fit\_transform} method was used on the training data and transforms it while learning distributional features from it for functions like scaling. The \texttt{transform} method applies the same transformations with the previously learned features and was used on the testing dataset. This sections lists all the transformations in the order they were applied.

\paragraph{Minimal profile flag}
A binary categorical feature was added, descerning between the accounts that have numeric features and the accounts without them. This step was done first because later steps add further numeric variables, which have missing values for different users, but the intended variable was intended to capture exactly the set of users that don't have numeric features in the initial dataset.

\paragraph{Subscribers Count Feature}
While reflecting a numeric quantity, the \texttt{subscribers\_count} feature was stored as a categorical variable with $1597$ possible categories. To enable the machine learning algorithms, to learn from the numeric relationships, this variable was converted to be numeric.

\paragraph{Data Scaling}
To ensure equal contribution of features and faster convergence for the logistic regressor, the data was scaled with the \texttt{StandarScaler} of the \texttt{sklearn} python library.

\paragraph{Median imputation}
The \texttt{SimpleImputer} class of the \texttt{sklearn} library was used to replace all missing values for numerical features with the median of this feature. That way, it is possible for the logistic regressor to use the numeric features without creating unnecessary noise. 

\paragraph{Categorical encoding}
Categorical features with exactly two categories were binary encoded, mapping each category to 0 or 1. The remaining categories were one-hot encoded, while assigning all categories with less than 30 members into an \texttt{other} bin. This \texttt{other} bin had the effect of reducing the number of city categories from 362 to only 6 and thereby preventing a drastic dimensionality increase that would have otherwise come from one-hot encoding it. For all categorical features, the \texttt{unknown} category was treated just as any other category.

\section{Models}

\subsection{Baseline}

For a baseline measurement, \texttt{sklearn}'s \texttt{DummyClassifier} was used with the "stratified" stragety. This strategy randomly predicts labels based on the frequency with which they occur in the dataset. It achieved an accuracy of 0.5 and an F1-score of 0.49 on the bot class.

\subsection{Logistic Regressor}

Logistic regression is one of the simplest machine learning algorithms. Despite its simplicity, it performs surprisingly well on linearly separable classification tasks. A major advantage of logistic regression is that its coefficients directly reflect feature importance for the classification task and can thus easily be interpreted. Additionally, it is very computationally efficient when dealing with large feature spaces. This makes it a good choice for local hyperparameter tuning and evaluation on the current dataset.

\subsection{XGBoost}

XGBoost (Extreme Gradient Boosting) is a high-performance implementation of gradient-boosted decision trees designed for speed and efficiency. \parencite{chenXGBoostScalableTree2016} It is particularly well-suited for structured datasets with a mix of categorical and numerical features, as in the current case.

One of XGBoost’s primary strengths lies in its ability to model complex, non-linear relationships between features and the target variable. Unlike logistic regression, which assumes linear separability, XGBoost can capture intricate interactions through its ensemble of decision trees. It also includes built-in regularization (both L1 and L2), which helps prevent overfitting—an important consideration when working with high-dimensional data.

Another advantage is its native handling of missing values. Instead of requiring imputation, XGBoost can learn optimal default directions for missing data during tree construction, making it more robust to incomplete feature sets. Furthermore, it supports parallelized training and early stopping, allowing for efficient model tuning even on larger datasets.


\section{Conclusion}
% Conclude with your main findings and possible future improvements.

Test

\printbibliography

\onecolumn
\appendix
\newpage

\section{Tables}
\renewcommand{\thetable}{A.\arabic{table}}
\setcounter{table}{0}

\begin{table}[H]
    \centering
    \caption{Sorted correlations between rings and other features}
    \label{tab:correlations}
    \begin{tabular}{lr}
    \toprule
     & Rings \\
    \midrule
 Shell weight & 0.628 \\
 Height & 0.610 \\
 Diameter & 0.575 \\
 Length & 0.557 \\
 Whole weight & 0.541 \\
 Viscera weight & 0.504 \\
 Shucked weight & 0.421 \\
    \bottomrule
    \end{tabular}
\end{table}
    

\vfill

\begin{table}[H]
    \centering
    \caption{Statistical comparisons of number of rings between genders}
    \label{tab:gender_comparisons}
    \begin{tabular}{llll}
        \toprule
        & M vs F & M vs I & F vs I \\
        \midrule
 t-statistic & -3.681 & 27.188 & 29.465 \\
 p-value & \textless{}0.001 & \textless{}0.001 & \textless{}0.001 \\
 Cohen's d & -0.139 & 1.016 & 1.154 \\
        \bottomrule
    \end{tabular}
\end{table}

\vfill

\newpage
\section{Figures}
\renewcommand{\thefigure}{A.\arabic{figure}}
\setcounter{figure}{0}

Test

\end{document}